
\documentclass{beamer}
\setbeamertemplate{navigation symbols}{}
\usepackage[catalan]{babel}
\usepackage[utf8]{inputenc}
\usepackage{ragged2e}
\usepackage{tikz}
\usepackage{amsmath}
\usepackage{amssymb}
\usepackage{stackrel}
\usetheme{Warsaw}
\justifying

\beamertemplatetransparentcovereddynamic
%\beamersetuncovermixins{\opaqueness<1>{25}}{\opaqueness<2->{15}}

\begin{document}
\title{Matemàtiques Segon Batxillerat}
\author{Artur Arroyo}
\date{\today}


\begin{frame}
\titlepage
\end{frame}

\begin{frame}\frametitle{Matemàtiques segon batxillerat}\tableofcontents[pausesubsections]
\end{frame}


\section{Matrius}

\begin{frame}{Definició de matriu}
\begin{block}{Definició}
Una matriu de \emph{m files i n columnes} és una taula de $m\times n$ nombres \emph{ordenats} de la forma:
\[A=\left(\begin{array}{cccccc}
a_{11} & a_{12} & a_{13} & \cdots & a_{1n}\\
a_{21} & a_{22} & a_{23} & \cdots & a_{2n}\\
a_{31} & a_{32} & a_{33} & \cdots & a_{3n}\\
\cdots & \cdots & \cdots & \cdots & \cdots\\
a_{m1} & a_{m2} & a_{m3} & \cdots & a_{mn}
\end{array}\right)
\]
El subíndex $i$ indica la fila i el $j$ la columna que ocupa cada element $a_{ij}$ de la matriu i sovint, es pot representar la matriu per $(a_{ij})$.
\end{block}
\end{frame}

\begin{frame}
\begin{block}{Dimensió d'una matriu}
Donada una matriu de  \emph{m files i n columnes} direm que la matriu és de dimensió $m\times n$.
\end{block}
\begin{exampleblock}{Exemple}
La matriu \[M=\left(\begin{array}{ccc}
-7 & \frac{3}{5} & 0\\
11 & -4 & 1
\end{array}\right)
\]
té dimensió $2\times3$, té dues files i tres columnes.
\end{exampleblock}
%\end{frame}
%\begin{frame}
\begin{block}{Matrius iguals}
Direm que dues matrius $A$ i $B$ són iguals si i només si tenen la mateixa dimensió i, a més, els elements coincideixen terme a terme. Dit d'una altra manera
$$A=B\Leftrightarrow a_{ij}=b_{ij}\quad\forall i,j$$
\end{block}
\end{frame}

\begin{frame}{Classificació de matrius}
\begin{itemize}
\item
Una matriu fila és aquella que té com a dimensió $1\times n$.
\item
Una matriu columna és una matriu que té com a dimensió $m\times1$.
\item
Una matriu zero és aquella en la qual tots els elements són zero.
\item
Una matriu quadrada és aquella que té el mateix nombre de files que de columnes. La seva dimensió és $n\times n$ i direm que la matriu és d'ordre $n$.
\end{itemize}
\end{frame}
\begin{frame}
\begin{exampleblock}{Exemples}
Siguin:
$A=\left(\begin{array}{ccc}
-2 & 7 & 0
\end{array}\right)
$
$B=\left(\begin{array}{c}
-4\\
2
\end{array}\right)
$
$C=\left(\begin{array}{cc}
0 & 0\\
0 & 0\\
0 & 0
\end{array}\right)
$\\ Llavors, $A$ és una matriu fila, $B$ és una matriu columna i $C$ és una matriu nul·la o zero de dimensió $3\times 2$.
\end{exampleblock}
\end{frame}

\begin{frame}{Tipus de matrius quadrades}
\begin{block}{Diagonal principal}
La diagonal principal d'una matriu està formada pels elements $a_{ii}$ de la matriu, és a dir, el $a_{11}$, el $a_{22}$, el $a_{33}$, etc.
\end{block}
Segons els elements que la formen, una matriu pot ser:
\begin{itemize}
\item
Matriu triangular superior: tots els elements situats per sota de la diagonal principal són zero.
\item
Matriu triangular inferior: tots els elements situats per sobre de la diagonal principal són zero.
\item
Matriu diagonal: tots els elements de la matriu són zero llevat dels de la diagonal principal.
\item
Matriu identitat o unitat: és una matriu diagonal en la que a més els elements de la diagonal principal són 1.
\end{itemize}
\end{frame}
\begin{frame}
\begin{exampleblock}{Exemples}
$A=\left(\begin{array}{ccc}
2 & 0 & 0\\
0 & -7 & 0\\
0 & 0 & 1\end{array}\right)$
$B=\left(\begin{array}{cc}
1 & 0\\
0 & 1\\
\end{array}\right)$
$C=\left(\begin{array}{cccc}
-3 & 0 & 0 & 0\\
2 & 4 & 0 & 0\\
7 & 5 & 1 & 0\\
4 & 6 & 9 & 6\end{array}\right)$

\bigskip
Llavors, $A$ és una matriu diagonal, $B$ és la matriu identitat d'ordre 2, que també s'escriu ($Id_{2}$), i $C$ és una matriu triangular inferior.
\end{exampleblock}
\end{frame}

\subsection{Operacions amb matrius}
\begin{frame}\frametitle{Suma de matrius}
\begin{block}{Definició}
La suma de dues matrius $A=(a_{ij})$ i $B=(b_{ij})$ dóna com a resultat una altra matriu $C$ que té per elements $c_{ij}=a_{ij}+b_{ij}$.
\end{block}
\begin{exampleblock}{Exemple}
Donades $A=\left(\begin{array}{ccc}
3 & 3 & 2\\
7 & -5 & 6\end{array}\right)$ $B=\left(\begin{array}{ccc}
4 & 5 & 0\\
-2 & 6 & -3\end{array}\right)$

llavors $A+B=\left(\begin{array}{ccc}
7 & 8 & 2\\
5 & 1 & 3\end{array}\right)$
\end{exampleblock}
\begin{alertblock}{Atenció!}
No és possible sumar dues matrius de dimensions diferents.
\end{alertblock}
\end{frame}
\begin{frame}\frametitle{Propietats de la suma de matrius}
Les propietas que té l'operació suma de matrius són
\begin{enumerate}
\item
Commutativa: $A+B=B+A$
\item
Associativa: $A+(B+C)=(A+B)+C$
\item
Element neutre: l'element neutre respecte la suma de matrius és la matriu zero de la dimensió corresponent. $A+0=A$
\item
Element oposat: l'element oposat (o invers respecte la suma) d'una matriu $A$ és la matriu que té tots els elements iguals que $A$ però amb el signe canviat. La matriu oposada de $A$ s'escriu $-A$.
\end{enumerate}
Aquesta darrera propietat ens permet, en particular, restar dues matrius de la següent manera $A-B=A+(-B)$.
\end{frame}

\begin{frame}\frametitle{Producte d'una matriu per un nombre real}
\begin{block}{Definició}
El producte d'una matriu $A$ per un nombre real $\alpha$ (també anomenat \emph{escalar}), dóna com a resultat una matriu $B$ de la mateixa dimensió que $A$ amb elements $b_{ij}=\alpha\cdot a_{ij}$.
\end{block}
\begin{exampleblock}{Exemple}
Donada la matriu $A=\left(\begin{array}{ccc}
2 & -3 & 5\\
0 & 1 & 7\end{array}\right)$, i $\alpha_{1}=2,$$\;\alpha_{2}=-3$ tenim $\alpha_{1}\cdot A=2\cdot\left(\begin{array}{ccc}
2 & -3 & 5\\
0 & 1 & 7\end{array}\right)=\left(\begin{array}{ccc}
4 & -6 & 10\\
0 & 2 & 14\end{array}\right)$, i $\alpha_{2}\cdot A=-3\cdot\left(\begin{array}{ccc}
2 & -3 & 5\\
0 & 1 & 7\end{array}\right)=
\left(\begin{array}{ccc}
-6 & 9 & -15\\
0 & -3 & -21\end{array}\right)$
\end{exampleblock}
\end{frame}
\begin{frame}{Propietats del producte d'un escalar per una matriu}
\begin{enumerate}
\item
Distributiva respecte la suma de matrius: $\alpha\cdot(A+B)=\alpha\cdot A+\alpha\cdot B$
\item
Distributiva respecte la suma d'escalars: $(\alpha+\beta)\cdot A=\alpha\cdot A+\beta\cdot A$
\end{enumerate}
\begin{alertblock}{Les matrius formen un $\mathbb{R}-$espai vectorial}
És clar, ja que si considerem el conjunt de les matrius de dimensió $m\times n$, fixada, $\mathcal{M}_{m\times n}$ tenim:
\begin{itemize}
\item
$\forall A, B\in \mathcal{M}_{m\times n},\quad A+B\in \mathcal{M}_{m\times n}$
\item
$\forall A\in \mathcal{M}_{m\times n},\quad \forall\alpha\in\mathbb{R},\quad \alpha\cdot A\in\mathcal{M}_{m\times n}$
\end{itemize}
\end{alertblock}
\end{frame}
\begin{frame}\frametitle{Producte de matrius}
\begin{block}{Definició}
El producte d'una matriu $A$, de dimensió $m\times n$ per una matriu $B$, de dimensió $n\times p$, és una altra matriu $C$ de dimensió $m\times p$, l'element $c_{ij}$ de la qual l'obtenim multiplicant la fila i-èssima de la primera matriu per la columna j-èssima de la segona matriu.
\end{block}
\begin{exampleblock}{Exemple}
Donades $A=\left(\begin{array}{ccc}
5 & -3 & 4\\
0 & 1 & 2\end{array}\right)$ i $B=\left(\begin{array}{cc}
4 & 2\\
0 & 5\\
1 & 3\end{array}\right)$

veiem que es poden multiplicar i que el seu producte serà una matriu de dimensió $2\times2$.
\end{exampleblock}
\end{frame}
\begin{frame}
\begin{exampleblock}{Producte de matrius}
 $$\stackbin[{\color{blue}2}\times{\color{red}3}]{}{\left(\begin{array}{ccc}
5 & -3 & 4\\
0 & 1 & 2\end{array}\right)}\cdot\stackbin[{\color{red}3}\times{\color{blue}2}]{}{\left(\begin{array}{cc}
4 & 2\\
0 & 5\\
1 & 3\end{array}\right)}=\stackbin[{\color{blue}2}\times{\color{blue}2}]{}{\left(\begin{array}{cc}
c_{11} & c_{12}\\
c_{21} & c_{22}\end{array}\right)}$$

Per calcular els $c_{ij}$ ho farem pas a pas. El terme $c_{11}$ ve de multiplicar la primera fila de la primera matriu per la primera columna de la segona matriu, com si es tractés
del producte escalar de dos vectors
$$\left(\begin{array}{ccc}
{\color{red}5} & {\color{orange}-3} & {\color{blue}4}\\
0 & 1 & 2\end{array}\right)\cdot\left(\begin{array}{cc}
{\color{red}4} & 2\\
{\color{orange}0} & 5\\
{\color{blue}1} & 3\end{array}\right)=\left(\begin{array}{cc}
{\color{red}5}\cdot{\color{red}4}+ ({\color{orange}-3})\cdot{\color{orange}0}+{\color{blue}4}\cdot{\color{blue}1}& c_{12}\\
c_{21} & c_{22}\end{array}\right)$$
\end{exampleblock}
\end{frame}
\begin{frame}
\begin{exampleblock}{Producte de matrius}
de forma que tenim $$\left(\begin{array}{ccc}
{\color{red}5} & {\color{red}-3} & {\color{red}4}\\
0 & 1 & 2\end{array}\right)\left(\begin{array}{cc}
{\color{red}4} & 2\\
{\color{red}0} & 5\\
{\color{red}1} & 3\end{array}\right)\left(\begin{array}{cc}
{\color{red}24} & c_{12}\\
c_{21} & c_{22}\end{array}\right)$$ i de forma similar, $$\left(\begin{array}{ccc}
{\color{red}5} & {\color{red}-3} & {\color{red}4}\\
0 & 1 & 2\end{array}\right)\left(\begin{array}{cc}
4 & {\color{red}2}\\
0 & {\color{red}5}\\
1 & {\color{red}3}\end{array}\right)\left(\begin{array}{cc}
24 & {\color{red}7}\\
c_{21} & c_{22}\end{array}\right)$$
\end{exampleblock}
\end{frame}

\begin{frame}
\begin{exampleblock}{Producte de matrius}
$$\left(\begin{array}{ccc}
5 & -3 & 4\\
{\color{red}0} & {\color{red}1} & {\color{red}2}\end{array}\right)\left(\begin{array}{cc}
{\color{red}4} & 2\\
{\color{red}0} & 5\\
{\color{red}1} & 3\end{array}\right)\left(\begin{array}{cc}
24 & 7\\
{\color{red}2} & c_{22}\end{array}\right)$$

i finalment, $$\left(\begin{array}{ccc}
5 & -3 & 4\\
{\color{red}0} & {\color{red}1} & {\color{red}2}\end{array}\right)\left(\begin{array}{cc}
4 & {\color{red}2}\\
0 & {\color{red}5}\\
1 & {\color{red}3}\end{array}\right)\left(\begin{array}{cc}
24 & 7\\
2 & {\color{red}11}\end{array}\right)$$
\end{exampleblock}
\end{frame}
\begin{frame}\frametitle{Propietats del producte de matrius}
\begin{enumerate}
\item
Associativa: $(A\cdot B)\cdot C=A\cdot(B\cdot C)$
\item
Element neutre: si $A$ té dimensions $m\times n$, llavors $Id_{m}\cdot A=A$ i $A\cdot Id_{n}=A$
\item
Distributiva respecte el producte:
\begin{enumerate}
\item
per l'esquerra: $A\cdot(B+C)=A\cdot B+A\cdot C$
\item
per la dreta: $(B+C)\cdot A=B\cdot A+C\cdot A$
\end{enumerate}
\end{enumerate}
\begin{alertblock}{Atenció!}
El producte de matrius en general no és commutatiu. D'entrada, per poder multiplicar dues matrius s'ha de complir la condició sobre les dimensions. Per una altra banda fins i tot matrius quadrades no commuten entre elles.
\end{alertblock}
\end{frame}
\begin{frame}
\begin{exampleblock}{Exemple 1.}
Siguin $\stackbin[m\times n]{}{A}$, $\stackbin[n\times r]{}{B}$, $\stackbin[r\times r]{}{C}$ i $\stackbin[r\times r]{}{D}$, llavors:
 \begin{itemize}
 \item
 els productes $\stackbin[m\times n]{}{A}\cdot\stackbin[n\times r]{}{B}$, $\stackbin[n\times r]{}{B}\cdot\stackbin[r\times r]{}{C}$ i $\stackbin[r\times r]{}{C}\cdot\stackbin[r\times r]{}{D}$ estan ben definits.
\item
Els productes $\stackbin[n\times r]{}{B}\cdot\stackbin[m\times n]{}{A}$ i $\stackbin[r\times r]{}{C}\cdot\stackbin[n\times r]{}{B}$ no estan ben definits per qüestió de dimensions.
\end{itemize}
\end{exampleblock}
\begin{exampleblock}{Exemple 2.}
 Finalment, en general $\stackbin[r\times r]{}{C}\cdot\stackbin[r\times r]{}{D}\neq\stackbin[r\times r]{}{D}\cdot\stackbin[r\times r]{}{C}$

Per exemple
$\left(\begin{array}{cc}
1 & 2\\
0 & 1\end{array}\right)\cdot\left(\begin{array}{cc}
1 & 0\\
2 & 1\end{array}\right)=\left(\begin{array}{cc}
5 & 2\\
2 & 1\end{array}\right)$
\bigskip
mentre que\\
a l'inrevés $\quad\left(\begin{array}{cc}
1 & 0\\
2 & 1\end{array}\right)\cdot\left(\begin{array}{cc}
1 & 2\\
0 & 1\end{array}\right)=\left(\begin{array}{cc}
1 & 2\\
2 & 5\end{array}\right)$
\end{exampleblock}
\end{frame}
%\end{frame}
\subsection{Matriu transposada}
\begin{frame}\frametitle{Matriu transposada}
\begin{block}{Definició}
Donada $A$ matriu de dimensió $m\times n$, definim la seva transposada $A^{t}$, com la matriu de dimensió $n\times m$ que té per files les columnes de $A$.
\end{block}
\begin{exampleblock}{Exemple}
Amb $A=\stackbin[2\times3]{}{\left(\begin{array}{ccc}
3 & -3 & 3\\
0 & 1 & 2\end{array}\right)}$, llavors $\; A^{t}=\stackbin[3\times2]{}{\left(\begin{array}{cc}
3 & 0\\
-3 & 1\\
3 &2 \end{array}\right)}$
\end{exampleblock}
\end{frame}
\begin{frame}\frametitle{Propietats de la matriu transposada}
\begin{enumerate}
\item
 $(A^{t})^{t}=A$
\item
$(A+B)^{t}=A^{t}+B^{t}$
\item
$(\alpha\cdot A)^{t}=\alpha\cdot(A)^{t}$
\item
$(A\cdot B)^{t}=B^{t}\cdot A^{t}$
\end{enumerate}
Cal recordar especialment l'ordre en que apareixen les matrius en la darrera propietat. El canvi es deu a la no commutativitat de les matrius.
\end{frame}
\begin{frame}\frametitle{Matrius simètriques i antisimètriques}
\begin{block}{Definició}
\begin{itemize}
\item
Diem que una matriu quadrada és \emph{simètrica} si coincideix amb la seva transposada, $A=A^{t}$
\item
Diem que una matriu quadrada és \emph{antisimètrica} si la seva oposada coicideix amb la seva transposada $-A=A^{t}$
\end{itemize}
\end{block}
\begin{block}{Descomposició en matriu simétrica i antisimétrica}
Sigui $A$ matriu quadrada qualsevol, llavors sempre es pot trobar una descomposició de $A$ com a suma d'una matriu simètrica i una altra antisimètrica.
$$A=\frac{A+A^{t}}{2}+\frac{A-A^{t}}{2}$$
\end{block}
\end{frame}
\subsection{Rang d'una matriu}
\begin{frame}\frametitle{Combinacions lineals de files en una matriu}
\begin{block}{Definició}
Diem que una fila $F_{i}$ no nul·la \emph{depèn linealment} de les files $F_{j_{1}},\; F_{j_{2}}\ldots F_{j_{m}}$ si se satisfà la igualtat:
$$F_{i}=\alpha_{1}F_{j_{1}}+\alpha_{2}F_{j_{2}}+\cdots +\alpha_{m}F_{j_{m}}$$
Diem que una fila és \emph{linealment independent} si no depèn linealment de cap altra fila de la matriu.
\end{block}
\end{frame}

\begin{frame}
\begin{exampleblock}{Exemple}
Si considerem la matriu $$\left(\begin{array}{cccc}
1 & 2 & 3 & 4\\
2 & 4 & 6 & 8\\
-1 & 2 & 3 & 4\\
0 & 4 & 6 & 8\end{array}\right)$$ Es pot comprovar fàcilment que $F_{2}= 2F_{1}$, i per tant $F_{2}$ no és linealment independent. De la mateixa manera, $F_{4}= F_{1}+F_{3}$ de forma que $F_{4}$ tampoc és linealment independent, queden doncs, com a files linealment independents, $F_{1}$ i $F_{3}$.

\end{exampleblock}
\end{frame}

\begin{frame}\frametitle{Rang d'una matriu}
\begin{block}{Definició}
El rang d'una matriu $A$, rang(A) és el nombre de files o de columnes linealment independents que té la matriu.
\end{block}
Per determinar el rang d'una matriu donada, en general ho farem sempre per files, i farem servir el que es coneix com a mètode de Gauss, que consisteix en convertir la matriu inicial en una que sigui triangular superior, és a dir, que tingui zeros a sota de la diagonal principal. Aquest procés es coneix com \emph{triangular la matriu} i el rang de la matriu serà el nombre de files no nul·les que té la matriu triangulada.
\end{frame}

\begin{frame}\frametitle{Càlcul del rang d'una matriu}
Hi ha certes transformacions permeses, anomenades \emph{transformacions elementals}, que deixen invariant el rang d'una matriu. Aquestes són:
\begin{itemize}
\item
Intercanviar dues files entre elles.
\item
Multiplicar una fila per un escalar
\item
Substituir una determinada fila per una combinació lineal d'altres.
\end{itemize}
\begin{exampleblock}{Exemple 1.}
$\left(\begin{array}{ccc}
1 & 2 & -1\\
3 & 0 & 1\end{array}\right)\sim\left(\begin{array}{ccc}
{\color{orange}1} & 2 & -1\\
0 & {\color{orange}6} & -4\end{array}\right)$  I per tant el rang és 2. Hem substituit la segona fila, $F_{2}$ per la combinació lineal $3F_{1}-F_{2}$ i així queda triangulada la matriu. Noteu el símbol ($\mathbf{\sim}$) que es fa servir per indicar que, tot i que les matrius \textbf{no} són iguals, tenen el mateix rang
\end{exampleblock}
\end{frame}
\subsection{Matriu inversa. Mètode de Gauss-Jordan}

\begin{frame}
\begin{exampleblock}{Exemple 2}
$\left(\begin{array}{cccc}
1 & 4 & 5 & 3\\
3 & 0 & 1\end{array}\right)\sim$
\end{exampleblock}

\end{frame}

\begin{frame}

\end{frame}


\end{document}
